\documentclass[10pt,a4paper]{amsart}
\usepackage[margin=19mm]{geometry}
\usepackage{amsmath,amssymb,mathtools}
\usepackage[scr=rsfs,cal=euler]{mathalfa}
\usepackage{xfrac}
\usepackage[unicode,hidelinks,bookmarks=false]{hyperref}
\newcommand{\ie}{i.e.\ }
\newcommand{\cf}{cf.\ }
\newcommand{\resp}{respectively\ }
\newcommand{\Subsection}{\subsection}
\providecommand{\nobreakdash}{\nobreak\hbox{-}}
\setlength{\emergencystretch}{2em}
\multlinegap0pt
\usepackage{mathtools}
\usepackage{amssymb}
\usepackage{enumerate}
\newtheorem{proposition}{Proposition}
\newtheorem{lemma}{Lemma}
\DeclareMathOperator{\Aut}{Aut}
\title{Polynomial-time isomorphism test for solvable groups with abelian Sylow subgroups}
\author{Saveliy V. Skresanov}
\date{Source: arXiv:2605.26748v2 (CC BY 4.0)}
\begin{document}
\maketitle

\noindent\emph{Source and licence.} Polynomial-time isomorphism test for solvable groups with abelian Sylow subgroups. Version arXiv:2605.26748v2 (CC BY 4.0), distributed by arXiv under the Creative Commons Attribution 4.0 International licence, http://creativecommons.org/licenses/by/4.0/, as recorded in the arXiv licence metadata for this exact version and shown on the abstract page https://arxiv.org/abs/2605.26748v2. Full licence notice: \texttt{licenses/arxiv-2605.26748-CC-BY-4.0.txt}.\par\smallskip

\noindent\textit{Editorial scope.} Counted unit: the article's lemma lred (source0.tex lines 641--646). The article's complete original proof of it is delivered with it (source0.tex lines 647--718, one \texttt{proof} environment of 72 lines). No mathematics is written, repaired or re-derived: the statement and the proof are the article's own lines, in the article's own notation, and no retained line is substituted or re-flowed.  Retained uncounted as the source-local context the proof needs: lines 541--546; lines 600--604.  Nothing was omitted from any retained span, so the package carries no omission record.  No retained line contains a citation, so the wrapper carries no bibliography and no bibliography entry.  No retained line contains a figure environment, an image-inclusion command or drawing code, so no image is delivered and the package declares no asset dependency.  Editorial formatting only, in addition: an AMS \texttt{amsart} preamble that restates the article's own declarations and macros used by the retained lines, namely the environments \texttt{proposition}, \texttt{lemma} on separate counters, together with the standard packages those lines need.  The retained environments print in this document as Proposition 1, Lemma 1 in order of appearance, whereas the article prints the same items with the numbering of its own full document.  The complete native source of the article as distributed in the arXiv e-print for this exact version is bundled unchanged as \texttt{sources/201440/source0.tex}.
\begin{proposition}\label{ptop}
	Let \( H \) and \( A \) be groups given by their Cayley tables, and assume we are given \( P \leq \Aut(H) \) 
	by permutation generators. Suppose that \( A \) is an abelian \( p \)-group, and \( H \) is a \( p' \)-group.
	If we are given representations \( \alpha, \beta : H \to \Aut(A) \) of \( H \) on \( A \), then one can compute
	\( P_{\alpha \to \beta} = \{ \phi \in P \mid \alpha^\phi \sim \beta \} \) in polynomial time.
\end{proposition}

\begin{proposition}\label{pbot}
	Let \( H \) and \( A \) be groups given by their Cayley tables, and assume that \( A \) is abelian.
	If \( \alpha, \beta : H \to \Aut(A) \) are representations of \( H \) on \( A \), then one can compute
	\( \Aut(A,\alpha \sim \beta) = \{ \psi \in \Aut(A) \mid \alpha^\psi = \beta \} \) in polynomial time.
\end{proposition}

\begin{lemma}\label{lred}
	Let \( G \) be an A-group given by its Cayley table, and assume we are given a characteristic
	abelian \( p \)-subgroup \( A \) of~\( G \). Suppose that \( A \) is characteristically complemented by \( H \) in \( G \),
	and we are given \( H \) and permutation generators of \( \Aut(H) \).
	Then we can compute permutation generators of \( \Aut(G) \) in polynomial time.
\end{lemma}

\begin{proof}
	Let \( \phi \in \Aut(G) \) be arbitrary. Since \( A \) is characteristically complemented by \( H \), there exists
	\( x \in G \) such that \( H^\phi = H^{x^{-1}} \). In particular, for an inner automorphism \( \iota_x(g) = x^{-1}gx \), \( g \in G \),
	the composition \( \phi \circ \iota_x \) stabilizes \( H \). Define \( S = \{ \phi \in \Aut(G) \mid H^\phi = H \} \)
	and observe that
	\[ \Aut(G) = \bigcup_{x \in G} S \cdot \iota_x = \langle S \cup \{ \iota_x \mid x \in G \} \rangle, \]
	so it suffices to compute the generators of \( S \).

	An element \( \phi \in S \) is completely determined by restrictions \( \phi|_A \in \Aut(A) \)
	and \( \phi|_H \in \Aut(H) \), so there is an embedding of \( S \) into \( \Aut(A) \times \Aut(H) \).
	By identifying \( S \) with its image under this embedding, we may assume that \( S \) is a subgroup of \( \Aut(A) \times \Aut(H) \).
	We will obtain generators of \( S \) by studying its projections \( S_A \) and \( S_H \) to \( \Aut(A) \) and \( \Aut(H) \), respectively.

	Define a representation \( \alpha : H \to \Aut(A) \) of \( H \) on \( A \) by
	\( a^{\alpha(h)} = a^h \), where \( a \in A \), \( h \in H \).
	Let \( h \in H \) be an arbitrary \( p \)-element. Since \( A \) is normal in \( G \),
	the group \( \langle A, h \rangle \) is a \( p \)-subgroup of \( G \). It is abelian, as \( G \) is an A-group,
	so \( h \) acts trivially on \( A \). Let \( K \) denote the subgroup of \( H \) generated by all \( p \)-elements.
	It follows that \( K \) lies in the kernel of \( \alpha \), and we have a representation of \( \overline{H} = H/K \)
	on \( A \) defined by \( \overline{\alpha}(hK) = \alpha(h) \), \( h \in H \).
	Note that \( K \), \( \overline{H} \), \( \alpha \) and \( \overline{\alpha} \) are computable in polynomial time.

	Since \( K \) is a characteristic subgroup of \( H \), any \( \eta \in \Aut(H) \) induces an automorphism
	\( \overline{\eta} \in \Aut(\overline{H}) \). The map sending \( \eta \in \Aut(H) \) to \( \overline{\eta} \in \Aut(\overline{H}) \)
	is a homomorphism from a permutation group on \( H \) to a permutation group on \( \overline{H} \), and
	we can compute its kernel \( C \) in polynomial time. Let \( P \) denote the image of this homomorphism, i.e.
	\[ P = \{ \overline{\eta} \in \Aut(\overline{H}) \mid \eta \in \Aut(H) \}, \]
	and observe that its generators can be computed in polynomial time as images of generators of \( \Aut(H) \).

	For any \( \eta \in \Aut(H) \), the subgroup \( K \) lies in the kernel of \( \alpha^\eta \),
	so \( \overline{\alpha^\eta} \) can be naturally defined. Clearly \( \overline{\alpha}^{\overline{\eta}} = \overline{\alpha^\eta} \),
	and \( \alpha \sim \alpha^\eta \) if and only if \( \overline{\alpha} \sim \overline{\alpha}^{\overline{\eta}} \).

	Let \( \nu \in \Aut(A) \), \( \eta \in \Aut(H) \), and note that the map \( \phi : G \to G \) defined by
	\( \phi(ah) = \nu(a)\eta(h) \), \( a \in A \), \( h \in H \), lies in \( S \) if and only if
	for all \( a \in A \), \( h \in H \) we have \( \phi(a^h) = \phi(a)^{\phi(h)} \).
	By the construction of \( \phi \) and definition of \( \alpha \) this is equivalent to
	\( \nu(a^{\alpha(h)}) = \nu(a)^{\alpha(\eta(h))} \), \( a \in A \), \( h \in H \).
	After applying \( \nu^{-1} \) to both sides we yield \( \alpha(h) = \nu \cdot \alpha(\eta(h)) \cdot \nu^{-1} \), \( h \in H \).
	It follows that \( \eta \) lies in the projection of \( S \) to \( \Aut(H) \) if and only if \( \alpha \sim \alpha^{\eta^{-1}} \), in other words
	\[ S_H = \{ \eta \in \Aut(H) \mid \alpha \sim \alpha^{\eta^{-1}} \}. \]
	By the previous paragraph, \( \alpha \sim \alpha^{\eta^{-1}} \) is equivalent to \( \overline{\alpha} \sim \overline{\alpha}^{\overline{\eta}^{-1}} \), so
	\[ S_H = \{ \eta \in \Aut(H) \mid \overline{\alpha} \sim \overline{\alpha}^{\overline{\eta}^{-1}} \}. \]

	Note that \( K \) contains all Sylow \( p \)-subgroups of \( H \), so \( |\overline{H}| \) is not divisible by~\( p \).
	We can now apply Proposition~\ref{ptop} to \( A \), \( \overline{H} \) and \( \overline{\alpha} \)
	and compute
	\[ P_{\overline{\alpha} \to \overline{\alpha}} = \{ \overline{\eta} \in P \mid
	\overline{\alpha} \sim \overline{\alpha}^{\overline{\eta}^{-1}} \} \]
	in polynomial time.
	The group \( S_H \) is the full preimage of \( P_{\overline{\alpha} \to \overline{\alpha}} \)
	under the map \( \overline{\phantom{a}} : \Aut(H) \to \Aut(\overline{H}) \). Therefore we can compute the generators of \( S_H \)
	in polynomial time by computing some preimages of generators of \( P_{\overline{\alpha} \to \overline{\alpha}} \) and adding the generators of~\( C \)
	which were computed earlier. Let \( \eta_1, \dots, \eta_t \) denote the generators of \( S_H \).

	For \( \eta \in S_H \) define \( S_A(\eta) = \{ \nu \in \Aut(A) \mid (\nu, \eta) \in S \} \).
	Note that \( S_A(\eta) \) is a coset of \( S_A(1_H) \), where \( 1_H \) denotes the identity automorphism of \( H \).
	For every \( \eta \in S_H \) choose some \( \nu_\eta \in \Aut(A) \) such that \( (\nu_\eta, \eta) \in S \). Then
	\[ S = \bigcup_{\eta \in S_H} S_A(\eta) \cdot (1_A, \eta) = \bigcup_{\eta \in S_H} S_A(1_H) \cdot (\nu_\eta, \eta) =
	\langle S_A(1_H) \cup \{ (\nu_\eta, \eta) \mid \eta \in S_H \} \rangle. \]
	We claim that \( S = \langle S_A(1_H) \cup \{ (\nu_{\eta_i}, \eta_i) \mid i = 1, \dots, t \} \rangle \).
	It is clear that the right hand side lies in \( S \), so it suffices to show that \( \langle (\nu_{\eta_i}, \eta_i) \mid i = 1, \dots, t \rangle \)
	contains a complete set of representatives for cosets of \( S_A(1_H) \) in \( S \).
	Given \( \eta \in S_H \), we can express it as \( \eta = \eta_{i_1} \cdots \eta_{i_k} \) for some \( i_1, \dots, i_k \in \{ 1, \dots, t \} \).
	Then \( (\nu_{\eta_{i_1}}, \eta_{i_1}) \cdots (\nu_{\eta_{i_k}}, \eta_{i_k}) = (\nu', \eta) \in S \) and hence
	\( S_A(\eta) \cdot (1_A, \eta) = S_A(1_H) \cdot (\nu', \eta) \), as required.

	Now the generators of \( S_A(1_H) = \Aut(A, \alpha \sim \alpha) \) can be computed in polynomial time by Proposition~\ref{pbot}.
	For every \( \eta_i \), \( i = 1, \dots, t \), we can find \( \nu_{\eta_i} \in \Aut(A) \) such that \( (\nu_{\eta_i}, \eta_i) \in S \)
	in polynomial time by Proposition~\ref{pbot} applied to \( \alpha \) and~\( \alpha^{\eta_i^{-1}} \).
	Hence we can compute the generators of \( S \) and thus of \( \Aut(G) \) in polynomial time.
\end{proof}

\end{document}
